\section{Shared Experimental Details}
\label{app:setup}

Sections~\ref{sec:normalization}--\ref{sec:density} contain the experiments;
this appendix records shared resource and evaluation details.

\paragraph{Training setup and run reuse.}
Section~\ref{sec:normalization} specifies the student, domains, and four
teacher weight sets. The $64$-response update comprises $16$ micro-batches
of four responses, with equal domain quotas, one optimizer update per
fresh rollout batch, and a $4{,}096$-token cap.
The initial budget is at most $500$ updates per run, subject to a
throughput pilot on a $96$-GB learner GPU and $48$-GB rollout/teacher GPUs.
Exact model revisions, optimizer and decoding settings, data splits,
and evaluation versions remain to be frozen before execution.

A representative single-teacher setting is chosen before inspecting
results. PG and corrected teacher top-$64$ in the single-teacher and
joint settings give four configurations. The joint top-$64$ configuration
with domain-response averaging is also the reference for the three-way
normalization study, which adds two averaging rules. Thus six training
configurations cover the core comparisons, with checkpoints reused for
the update, teacher-overlap, and full-vocabulary measurements.
Additional teacher settings, the response-cap comparison, and repeated
seeds require separate resource accounting. Final training outcomes remain
pending; Section~\ref{sec:density} records preliminary observations.

\paragraph{Capability evaluation.}
Use MATH-500 greedy pass@1, a frozen LiveCodeBench pass@1 slice,
IFBench strict accuracy, and GPQA-Diamond average@4. Record evaluator,
model, and data revisions and fix decoding within each comparison.
Report raw per-domain scores, their macro-average, and the worst-domain
change from initialization. Independently trained OPD students, where
available, are references at matched domain exposure, not upper bounds
or guarantees of jointly attainable capability. A lack of single-teacher
transfer is not itself evidence of multi-teacher interference.
Loss on a common held-out response bank and loss on fresh student
responses are reported separately; a lower mean KL alone does not
establish balanced capability integration.

\paragraph{Teacher reference results.}
Table~\ref{tab:teacher_reference} reports the completed reference evaluations
from September 6, 2026. Each RL-trained teacher is evaluated only on its
own domain; the initial Qwen3-1.7B-Base student is evaluated on the same
questions and labels. Scores include truncated responses.

\begin{table}[htbp]
\centering
\small
\setlength{\tabcolsep}{4pt}
\begin{tabular}{lrrrrr}
\toprule
Benchmark & Responses & Teacher & Initial & $\Delta$ (pp) & Trunc. \\
\midrule
MATH-500 & 500 & 75.60 & 55.60 & +20.00 & 6.80 \\
LiveCodeBench postcutoff & 128 & 23.44 & 0.78 & +22.66 & 24.22 \\
IFBench strict & 300 & 29.00 & 17.33 & +11.67 & 7.33 \\
GPQA-Diamond avg@4 & $198\times4$ & 34.60 & 18.81 & +15.78 & 0.51 \\
\bottomrule
\end{tabular}
\caption{Domain-specific teacher references and the initial student.
Scores and teacher truncation rates (Trunc.) are percentages;
$\Delta$ is the teacher-minus-initial difference in percentage points,
computed before rounding. MATH-500 and LiveCodeBench use pass@1;
IFBench uses strict accuracy. GPQA averages correctness over four
responses per question, rather than counting any correct response as
success. Each row uses a different teacher and benchmark.}
\label{tab:teacher_reference}
\end{table}

The teachers produce respectively $378/500$, $30/128$, $87/300$, and
$274/792$ correct responses. Mathematics, code, and instruction following
use greedy decoding; GPQA uses temperature $1$ and top-$p=1$.
The requested response limit is $32{,}768$ tokens, reduced as needed to
fit each prompt within the $32{,}768$-token context.
Teacher prompts retain the empty \texttt{<think>}\ldots\texttt{</think>}
suffix from the chat template, whereas initial-student prompts remove
it. Thus the observed differences compare the deployed evaluation
configurations and do not isolate the effect of model weights.

Code generation warrants particular attention: $31/128$ teacher responses
are truncated and none passes; $27/128$ have execution errors, with zero
sandbox infrastructure errors. These categories can overlap.
The initial student's code and IF truncation rates are $75.78\%$ and
$69.67\%$, respectively, so generation behavior also matters when
interpreting their score differences. These reference scores do not
establish transfer to an OPD student or jointly attainable performance.
% Auditable counts, source paths, completion records, and artifact hashes:
% experiments/teacher_reference_results_20260906.json

\paragraph{Compute and uncertainty.}
Record generated tokens, domain response counts, optimizer updates, and
GPU hours separately, including teacher-scoring and diagnostic overhead.
Key online comparisons use repeated paired training seeds when feasible;
initial one-seed runs are exploratory. Prompt bootstrap estimates
conditional evaluation uncertainty, not training-seed variation.
Report effect sizes, confidence intervals, and the range of differences
the experiment can resolve. Teacher pairs sharing weights or training
lineage, and repeated checkpoints, are not independent observations.
No sparsity advantage or capability improvement is claimed before
measurement.
